\documentclass[11pt]{article}
\usepackage[a4paper,margin=27mm]{geometry}
\usepackage[T1]{fontenc}
\usepackage{lmodern,microtype,amsmath,amssymb,amsthm,mathtools,enumitem}
\usepackage[colorlinks=true,linkcolor=blue]{hyperref}
\newtheorem{theorem}{Theorem}
\newtheorem{lemma}[theorem]{Lemma}
\newtheorem{proposition}[theorem]{Proposition}
\theoremstyle{remark}
\newtheorem{remark}[theorem]{Remark}
\newcommand{\R}{\mathbb R}
\newcommand{\supp}{\operatorname{supp}}
\newcommand{\sd}{\mathbin{\triangle}}
\newcommand{\Pc}{\mathcal P}
\newcommand{\one}{\mathbf 1}
\newcommand{\F}{\mathcal F}
\title{An Ordinal Refinement by a Slack Envelope}
\author{}
\date{7 October 2026}
\begin{document}
\maketitle

\paragraph{Scope.}
This note proves a simultaneous refinement of the basic level extension in
\emph{Loss Relations from Credal Families}. The refinement respects strict
background inclusion and has closed competitive choice. It uses incomparability,
not a real-valued loss. Only continuity of the family is needed; neither a
Lipschitz bound, a domination bound on upper probabilities, nor a quantitative
estimate on error measures is required. The structural compatibility of covers
is still an assumption. The refinement theorem by itself therefore leaves open
the existence of suitable families through every credal set in every dimension.

The unrestricted refinement lemma based only on background inclusion,
admissibility, competitiveness and closed sections is false; a counterexample
appears at the end. The positive result below exploits the additional slack
structure of the basic level extension. The companion note
\emph{Full Coverage by Homothetic Credal Families} (8 October 2026) now
constructs families satisfying the assumptions below through every credal set,
including all vertex-containing targets, and thereby supplies full (A).

\section{Assumptions and the basic extension}

Fix $n\ge2$, the simplex $\Delta\subset\R^n$, and a finite probability measure
$\mu$ with the same null sets as Lebesgue measure (for example the standard
Gaussian). Models are all Lebesgue-measurable subsets of $\R^n$. Put
$G_i=\{g:g_i>0\}$ and $E_i(X)=X\sd G_i$. Write
\[
 X\le_iY\quad\Longleftrightarrow\quad
 E_i(X)\subseteq E_i(Y)\text{ a.e.},
 \qquad X\approx Y\quad\Longleftrightarrow\quad\mu(X\sd Y)=0.
\]
The strict part of $\le_i$ is $<_i$. Thus $X<_iY$ means a positive-measure
reduction of errors.

Let $p\mapsto K_p$ be a family of nonempty closed convex subsets of $\Delta$,
and let
\[
 D_p=\{g:\min_{m\in K_p}m\cdot g>0\}\cup\{g:g\ge0,\ g\ne0\}.
\]
Assume:
\begin{enumerate}[label=\textbf{(\arabic*)},leftmargin=*]
\item \emph{Face support:} $K_p\subseteq\Delta_{\supp p}$.
\item \emph{Continuity:} $\mu(D_p\sd D_q)\to0$ as $q\to p$.
\item \emph{Compatibility of covers:} for every model $X$, $i\ne j$, and
indices $p,q$,
\[
 X\le_iD_p,\quad X\le_jD_q\quad\Longrightarrow\quad p_i\le q_i.
 \tag{CC}\label{eq:cc}
\]
\end{enumerate}
Condition (CC) follows from (U) in the credal-family note. It also follows from
the proposed two-sided (U$^*$): the common-cover gluing argument supplies
$\overline U_p(i)\le\overline U_q(i)$ and
$\overline U_q(j)\le\overline U_p(j)$, and (U$^*$) yields $p_i\le q_i$.
Hausdorff continuity of $K_p$ implies condition (2): the lower expectations
converge pointwise, their zero sets are null, and dominated convergence applies
to acceptance indicators.

Define
\[
 \Pc_i(X)=\{p:X\le_iD_p\},\qquad
 T(X)=\{i:\Pc_i(X)\ne\varnothing\}.
\]
The sets $\Pc_i(X)$ are closed: the function
$p\mapsto\mu(E_i(X)\setminus E_i(D_p))$ is continuous by (2).
For $i\in T(X)$ put
\[
 a_i(X)=\max_{p\in\Pc_i(X)}p_i.
\]
Two extensions of this notation will be useful. Set
\[
 \bar a_i(X)=\begin{cases}a_i(X)&i\in T(X),\\0&i\notin T(X),\end{cases}
 \qquad
 \alpha_i(X)=\begin{cases}a_i(X)&i\in T(X),\\-1&i\notin T(X).\end{cases}
\]
The first is used for budgets, the second for comparisons. Let
\[
 s(X)=1-\sum_i\bar a_i(X),\qquad
 \F=\{X:X\approx D_p\text{ for some }p\in\Delta\}.
\]
Under these assumptions the basic extension with level orders is precisely
\[
 X\preceq_i^*Y\quad\Longleftrightarrow\quad
 \alpha_i(X)\ge\alpha_i(Y).
 \tag{*}\label{eq:basicorder}
\]
Indeed, a nonempty cover set has an attained maximum $i$th coordinate, so its
upward closure in the level order is determined by that maximum. Empty cover
sets give the common bottom level $-1$.

\section{Compactness at zero slack}

\begin{lemma}\label{lem:elementary}
The following hold.
\begin{enumerate}[label=(\alph*)]
\item $X\le_iY$ and $i\in T(Y)$ imply $i\in T(X)$ and
$a_i(X)\ge a_i(Y)$.
\item $a_i(D_p)=p_i$ for every $p,i$.
\item If $D_p\le_iD_q$ and $p_i=q_i$, then $p=q$.
\item $s(X)\ge0$, and $s(X)=0$ if and only if $X\in\F$.
\end{enumerate}
\end{lemma}
\begin{proof}
(a) Cover sets are nested. For (b), $p$ is itself a cover, while any
$q\in\Pc_i(D_p)$ can be paired with $p\in\Pc_j(D_p)$ for $j\ne i$ in (CC),
giving $q_i\le p_i$. For (c), for each $j\ne i$ apply (CC) to
$p\in\Pc_j(D_p)$ and $q\in\Pc_i(D_p)$ to obtain $p_j\le q_j$.
Equality at $i$ and equal coordinate sums force $p=q$.

For (d), suppose $T=T(X)\ne\varnothing$, and choose maximizing covers
$p^j\in\Pc_j(X)$, $j\in T$. For $k\in T$, (CC) gives
$p^j_k\ge p^k_k=\bar a_k(X)$, with equality when $k=j$; for $k\notin T$ the
same inequality follows from nonnegativity. Hence
\[
 p^j\ge\bar a(X),\qquad
 \|p^j-\bar a(X)\|_1=1-\sum_k\bar a_k(X)=s(X)\ge0.
 \tag{1}\label{eq:distance}
\]
If $s(X)=0$, every maximizing cover equals $p=\bar a(X)\in\Delta$, and
$\supp p\subseteq T$. Thus $X\le_iD_p$ for every $i\in\supp p$.
Face support implies $X\approx D_p$: an added gamble would have to be positive
at every world in $\supp p$, and hence already accepted by $D_p$; a removed
gamble would have to be nonpositive at every such world, and hence not accepted
by $D_p$, apart from a null set on coordinate hyperplanes.
Conversely $X\approx D_p$ gives $s(X)=0$ by (b). If $T$ is empty, $s(X)=1$.
\end{proof}

\begin{lemma}[Stability at zero slack]\label{lem:stability}
If $s(X_k)\to0$, then some subsequence admits $p\in\Delta$ such that
\[
 \bar a(X_k)\longrightarrow p,
 \qquad \mu(X_k\sd D_p)\longrightarrow0.
\]
\end{lemma}
\begin{proof}
Pass to a subsequence with the nonempty set $T(X_k)=T$ fixed and
$\bar a(X_k)\to p$. The coordinate sums tend to one, so $p\in\Delta_T$.
For $j\in T$ choose maximizing covers $p^{j,k}$. Equation
\eqref{eq:distance} gives $p^{j,k}\to p$, and (2) gives
$\mu(D_{p^{j,k}}\sd D_p)\to0$.

Up to a null set, if $X_k$ differs from $D_p$ at a gamble $g$, some
$j\in T$ regards that difference as an extra error relative to $D_p$.
Indeed, if $g\in X_k\setminus D_p$, then $g_j\le0$ for some $j\in T$,
since all-positive payoffs on $T$ force acceptance by $D_p$. If
$g\in D_p\setminus X_k$, then $g_j>0$ for some $j\in T$, since otherwise
face support forces rejection by $D_p$, except on coordinate hyperplanes.
As $X_k\le_jD_{p^{j,k}}$, the corresponding covering model must differ from
$D_p$ at $g$. Therefore
\[
 \mu(X_k\sd D_p)
 \le\sum_{j\in T}\mu(D_{p^{j,k}}\sd D_p)\longrightarrow0.
\]
\end{proof}

\section{The envelope}

For $i\in T(X)$ define
\[
 b_i(X)=\inf_{Y\le_iX}
       \left(a_i(Y)+\frac{s(Y)}{2n}\right).
 \tag{2}\label{eq:envelope}
\]
Every $Y$ in this infimum is covered at $i$ by Lemma~\ref{lem:elementary}(a).
The infimum is over a nonempty set, since it includes $Y=X$.

\begin{lemma}[Positive monotone reserve]\label{lem:envelope}
For covered models:
\begin{enumerate}[label=(\alph*)]
\item $a_i(X)\le b_i(X)\le a_i(X)+s(X)/(2n)$;
\item $X\le_iY$ implies $b_i(X)\ge b_i(Y)$ whenever $Y$ is covered at $i$;
\item $b_i(X)=a_i(X)$ if and only if $X\in\F$.
\end{enumerate}
In particular,
\[
 \sum_{i\in T(X)}b_i(X)\le1-\frac{s(X)}2.
 \tag{3}\label{eq:newbudget}
\]
\end{lemma}
\begin{proof}
(a) Every term in the infimum is at least $a_i(X)$ by monotonicity of $a_i$;
the term $Y=X$ gives the upper bound. For (b), the background lower set of
$X$ is contained in that of $Y$, so its infimum is no smaller.

If $X\in\F$, (a) and $s(X)=0$ give equality. Conversely suppose
$b_i(X)=a_i(X)$, and choose $Y_k\le_iX$ with
\[
 a_i(Y_k)+\frac{s(Y_k)}{2n}\longrightarrow a_i(X).
\]
Since $a_i(Y_k)\ge a_i(X)$ and $s(Y_k)\ge0$, we have
$a_i(Y_k)\to a_i(X)$ and $s(Y_k)\to0$. Lemma~\ref{lem:stability} gives,
along a subsequence, $Y_k\to D_p$ in symmetric-difference measure, with
$p_i=a_i(X)$. Passing to the limit in $Y_k\le_iX$ gives $D_p\le_iX$.
Choose a maximizing cover $q\in\Pc_i(X)$, so $q_i=a_i(X)=p_i$. Then
\[
 D_p\le_iX\le_iD_q.
\]
Lemma~\ref{lem:elementary}(c) gives $p=q$, and the displayed sandwich gives
$X\approx D_p$. This proves (c).

Finally sum (a):
$\sum_{i\in T(X)}b_i(X)\le1-s(X)+|T(X)|s(X)/(2n)\le1-s(X)/2$.
\end{proof}

\paragraph{Why this addresses the induction problem.}
The lift at $X$ is
\[
 b_i(X)-a_i(X)
 =\inf_{Y\le_iX}\left(a_i(Y)-a_i(X)+\frac{s(Y)}{2n}\right).
\]
Thus it accounts simultaneously for the gap to every background improvement
$Y$ and for that model's reserve. The infimum could a priori be zero.
Lemma~\ref{lem:envelope} proves it is positive for every covered nonfamily
model. The proof is qualitative compactness, not a uniform numerical lower
bound on all lifts. There is no accumulation of successive expenditures:
all reserves are defined at once, with the common bound \eqref{eq:newbudget}.

\section{The partial ranking}

Put
\[
 v_i(X)=
 \begin{cases}(a_i(X),b_i(X))&i\in T(X),\\(-1,-1)&i\notin T(X).
 \end{cases}
\]
Larger endpoints mean better performance. Define $X\preceq_iY$ when
$v_i(X)\ge v_i(Y)$ coordinatewise and at least one of the following holds:
\begin{enumerate}[label=(\roman*)]
\item $v_i(X)\ne v_i(Y)$;
\item the common value is $(t,t)$ for some $t\ge0$;
\item $X\le_iY$.
\end{enumerate}
Thus different endpoint pairs are compared by the product order. Within an
equal-pair class, retain background inclusion, except that all models at a
nonnegative degenerate pair tie. By Lemma~\ref{lem:envelope}, those exceptional
classes consist entirely of family models modulo null sets.

\begin{theorem}\label{thm:refinement}
Under assumptions (1)--(3), these relations are preorders and form a
subrelation of the basic extension \eqref{eq:basicorder}. They satisfy:
\begin{enumerate}[label=(\alph*)]
\item $X\le_iY\Rightarrow X\preceq_iY$, and
$X<_iY\Rightarrow X\prec_iY$;
\item for every $q$, $D_q$ is competitive on $\supp q$;
\item every set $\{q:D_q\preceq_iX\}$ is closed;
\item if $X\preceq_iD_q$ for every $i$, then $X\approx D_q$;
\item every $X\notin\F$ is strictly dominated at every world by some $D_q$.
\end{enumerate}
Consequently (E), (LA), (S), and (C) hold, with all measurable models
evaluable and coherence understood modulo null sets. The undominated models
are exactly $\F$. The construction does not assert family-wide (A) unless
the families under consideration collectively contain every intended coherent
model modulo null sets.
\end{theorem}
\begin{proof}
Reflexivity is immediate. For transitivity, endpoint pairs decrease
coordinatewise along a chain. If the first and last pairs differ, clause (i)
applies. If they agree, all pairs agree; the comparison inside that class is
either universal or background inclusion, both transitive. The first endpoint
is $\alpha_i$, so the new relation is a subrelation of \eqref{eq:basicorder}.

(a) Both endpoints are background-monotone, including transitions from a
covered model to an uncovered one. A background comparison with equal endpoint
pairs is retained by (iii). If the background comparison is strict, unequal
pairs exclude its reverse. With equal pairs, the only possible exception would
be clause (ii). But then $X\approx D_p$, $Y\approx D_q$, and $p_i=q_i$.
Background inclusion gives $D_p\le_iD_q$, whence $p=q$ by
Lemma~\ref{lem:elementary}(c). Thus $X\approx Y$, contradicting strictness.
In all other equal-pair classes the strict background comparison remains strict.

For (b) and (c), the exact selected-model comparisons are
\[
 D_q\preceq_iX\quad\Longleftrightarrow\quad
 \begin{cases}q_i\ge b_i(X)&i\in T(X),\\\text{always}&i\notin T(X),\end{cases}
 \tag{4}\label{eq:outgoing}
\]
and
\[
 X\preceq_iD_q\quad\Longleftrightarrow\quad
 i\in T(X)\ \text{and}\ a_i(X)\ge q_i.
 \tag{5}\label{eq:incoming}
\]
Indeed $v_i(D_q)=(q_i,q_i)$. If $X\notin\F$ is covered, its endpoints are
strictly separated, so at either threshold the product comparison is strict
in one endpoint. If $X\in\F$, equal endpoints tie by (ii).

If some $i\in\supp q$ is uncovered, competitiveness is immediate. Otherwise,
if $q_i<b_i(X)$ held for every $i\in\supp q$, then
$1=\sum_{i\in\supp q}q_i<\sum_{i\in\supp q}b_i(X)\le1$, a contradiction
to \eqref{eq:newbudget}. Equation~\eqref{eq:outgoing} also proves closedness.

(d) Equation~\eqref{eq:incoming} gives $T(X)=[n]$ and $a_i(X)\ge q_i$
for all $i$. The budget forces equality at every coordinate and $s(X)=0$.
The proof of Lemma~\ref{lem:elementary}(d) then yields $X\approx D_q$.

(e) Put $\beta_i=b_i(X)$ for $i\in T(X)$ and $\beta_i=0$ otherwise.
Since $X\notin\F$, $s(X)>0$, and \eqref{eq:newbudget} gives
$\sum_i\beta_i<1$. Let
\[
 q=\beta+\frac{1-\sum_i\beta_i}{n}\one.
\]
For a covered world $q_i>b_i(X)\ge a_i(X)$, and for an uncovered world the
family endpoint pair is strictly above $(-1,-1)$. In either case
$D_q\prec_iX$.

Finally, all null modifications of a family model have its profile. Part (d)
rules out domination of family models, while (e) rules out every other
undominated model. All family models are coherent. These observations establish
the stated axioms and characterize the undominated models.
\end{proof}

\begin{remark}[Where incomparability does the work]
Let $X\notin\F$ be covered at $i$, with endpoints $a_i(X)<b_i(X)$.
For every family index $r$ with
\[
 a_i(X)<r_i<b_i(X),
\]
$D_r$ and $X$ are incomparable at $i$. At $r_i=a_i(X)$, $X$ is strictly
better; at $r_i=b_i(X)$, $D_r$ is strictly better. Thus the comparison section
is the closed set $\{r:r_i\ge b_i(X)\}$.
The refinement deletes some strict comparisons of $\preceq^*$ without reversing
them. It therefore lies outside the hypothesis of the note's no-selection
proposition, which requires preservation of every strict comparison of
$\preceq^*$. Endpoint pairs are not real-valued losses: the resulting preorder
has genuine incomparabilities.
\end{remark}

\section{Why the unrestricted refinement lemma is false}

\begin{proposition}\label{prop:counterexample}
Background inclusion, undominated selected coherent models, competitiveness,
and closed comparison sections do not alone guarantee that a strict background
improvement can be made strict by passing to a subrelation while preserving
the same closed competitive choice.
\end{proposition}
\begin{proof}
Take two worlds, all measurable models, the unit ball $B\subset\R^2$, and
the precise family
\[
 D_p=\{g:p\cdot g>0\}\cup\{g:g\ge0,\ g\ne0\}.
\]
For this counterexample only, let the initial rankings be generated by
\[
 L_i(X)=\int_B |g_i|\,1_{X\sd G_i}(g)\,dg.
\]
They respect weak background inclusion. For each prior $p$, the unique Bayes
acceptance region inside $B$, modulo null sets, is $D_p\cap B$, since the
expected advantage of acceptance is $p\cdot g$ and its zero set is null.
Therefore $D_p$ is undominated, is competitive on $\supp p$, and has losses
continuous in $p$. Distinct indices give distinct acceptance regions inside
$B$, and no family member weakly dominates a different family member.

Fix $q=(1/2,1/2)$ and a nonempty open set
$A\subset B^c\cap G_1\cap\{g:q\cdot g<0\}$. Let $X=D_q\cup A$.
Then $X<_1D_q$, but $L(X)=L(D_q)$.

Suppose a subrelation retained preorders and closed competitiveness for the
same family while removing $D_q\preceq_1X$. The KKM argument applied to $X$
would give some $p$ with $D_p\preceq_iX$ at both worlds. Since these are
subrelations of the original rankings, $L(D_p)\le L(X)=L(D_q)$, forcing
$p=q$. But this requires the removed comparison $D_q\preceq_1X$, a
contradiction.
\end{proof}

This obstruction is absent from the positive theorem: zero slack identifies
the family itself modulo null sets, and Lemma~\ref{lem:stability} controls
sequences approaching zero slack. The positive theorem proves a successful
simultaneous refinement of the canonical basic extension; it does not claim
that every intermediate subrelation satisfying the weaker conditions admits
every further refinement.

\end{document}
